\documentclass{ltjsbook}
\usepackage{docmute} 
\usepackage{../../myStyle} 

\begin{document}
\VerbatimFootnotes
% \chapter{モデリングの目から見た検定3；多群の平均値差モデル}
\chapter{Test 3 from the Perspective of Modeling: Multiple Group Mean Difference Model}
\label{x21_modeling3}
% \section{要因計画モデル}
\section{Factorial Design Models}
% さてここまでは二水準の平均値差を考えるモデルでしたが，ここからそれを三水準に増やしてみましょう。
Until now, we've been considering a model with two levels of average difference. However, let's try increasing that to three levels from here.
% 水準数が多くなると\keyterm{要因計画}{よういんけいかく}{Factorial Design}や\keyterm{実験計画}{じっけんけいかく}{Experimental Design}といった話になってきます。\keyterm{分散分析}{ぶんさんぶんせき}{ANalysis Of VAriance;ANOVA}は要因計画と帰無仮説検定の合わせ技で，効果がある/ないの判定を分散の比をつかって行うことでした\footnote{忘れた人は，データ解析基礎のテキストに戻って確認してください。}。
When the number of levels increases, it comes to discussions on Factorial Design and Experimental Design. ANOVA (ANalysis Of VAriance) is a combination of factorial design and null hypothesis testing, where decisions are made on whether or not there are effects by using the ratio of variance. If you forgot, please go back and check the basics of data analysis in the text.

% 今回は要因計画をデータ生成モデルから考え，\keyterm{ベイズ推定}{べいずすいてい}{Bayesian Inference}をすることになります。\keytermL{帰無仮説検定}{きむかせつけんてい}のようなYes/No判断はせず，平均値や差の大きさを直接見積もることになるのは，二群の時と同じです。
This time, we will consider the factor plan from the data generation model and perform \keyterm{Bayesian Inference}. Like the \keytermL{null hypothesis test}, there will be no Yes/No decisions, and we will directly estimate the size of the mean or difference, just as we did when dealing with two groups.

% まずは簡単な\keyterm{群間計画}{ぐんかんけいかく}{Between Design}の話から進めましょう。対応のない二群の時のように，それが三群になるモデルですので，モデルの設計図は非常に単純です(図\ref{fig::21_01})。
Let's first proceed from the discussion about the simple \keyterm{Between Design}. Like the times when there is no correspondence between the two groups, it is a model that becomes three groups, so the design diagram of the model is very simple (Figure \ref{fig::21_01}).
\begin{figure}[ht]
  \begin{center}
    \includegraphics[width=0.8\linewidth,keepaspectratio]{../images/21_modeling3/Slide21_01.png}
    \caption{多群の平均値の差を比べるときの設計図}
    \label{fig::21_01}
  \end{center}
\end{figure}

% 設計図を見れば，コードはそのまま置き換える形でかけますね。設計図の通りに書いてみたコードがcode:\ref{stan::21_01}にあります。
Looking at the design drawing, the code can be written just by replacing it as it is. The code written according to the design drawing is in code:\ref{stan::21_01}.

\begin{lstlisting}[language=C++,label={stan::21_01},caption={三群の平均値のコード}]
data{
  int<lower=0> N1;
  int<lower=0> N2;
  int<lower=0> N3;
  array[N1] real X1;
  array[N2] real X2;
  array[N3] real X3;
}

parameters{
  real mu1;
  real mu2;
  real mu3;
  real<lower=0> sig;
}

model{
  // likelihood
  X1 ~ normal(mu1,sig);
  X2 ~ normal(mu2,sig);
  X3 ~ normal(mu3,sig);
  // prior
  mu1 ~ uniform(0,100);
  mu2 ~ uniform(0,100);
  mu3 ~ uniform(0,100);
  sig ~ cauchy(0,5);
}

generated quantities{
  real diff12;
  real diff13;
  real diff23;
  
  diff12 = mu1 - mu2;
  diff13 = mu1 - mu3;
  diff23 = mu2 - mu3;
}

\end{lstlisting}


% 最後の\verb|generated quantities|ブロックにあるように，各群の差の分布も計算して出せるようにしてあります。これをみながら，群AとB，AとC，BとCなどどこに差があるのか，その大きさはどれぐらいかということを考えることができます。左の大きさを標準偏差\verb|sig|で割ることで効果量を出すこともできます。
As in the last \verb|generated quantities| block, it is set up to be able to calculate and output the distribution of differences between each group. Looking at this, one can consider where there are differences between group A and B, A and C, B and C, and how large these differences are. The effect size can also be obtained by dividing the size on the left by the standard deviation \verb|sig|.

% 二群の平均値差の時もそうでしたが，検定と違って差の大きさを直接検証できているところがポイントです。また，検定の場合は分析によって有意差が確認できた場合，そのあと\keyterm{多重比較}{たじゅうひかく}{multiple comparison}へと進むのでした。これは帰無仮説が$H_0:\mu_1 = \mu_2 = \mu_3$となっていて，これを否定することから「どこに差があったか」をさらに詳しくみる必要があったからです。この時，統計的検定の繰り返しによって\keyterm{タイプ1エラー}{たいぷ1えらー}{Type I Error}が増えること，すなわち$\alpha = 0.05$と5\% 水準で検定していても，検定を繰り返すと5\% を大きく上回る問題が生じるため，有意水準を調節するなど工夫が必要だという話でした\footnote{お前は何をいってるんだ，と思った人は，データ解析基礎のテキストに戻って確認しておいてください。大雑把に解説しますと，5\%水準というのは一回の判断で間違いが含まれる可能性を5\%にしましょう，という判定基準のことでしたが，二回判定を繰り返すと$1-(0.95^2)=0.0975$と9.75\%水準になってしまうということでした。3群の平均値を比べるときは，3箇所比較する必要があり，\keytermL{t検定}{tけんてい}を3回も繰り返すと5\%水準ではなくなる，という問題です。}。
As it was the case with the difference in average values of two groups, the point is that you can directly verify the size of the difference unlike a test. Also, in the case of a test, if a significant difference can be confirmed by the analysis, it will proceed to \keyterm{multiple comparison}. This was because the null hypothesis was $H_0:\mu_1 = \mu_2 = \mu_3$, and it was necessary to look more in detail at "where the difference was" by denying it. In this case, \keyterm{Type I Error} increases by repeating statistical tests, that is, even if you test at an $\alpha = 0.05$ and a 5% level, if you repeat the test, a problem arises that greatly exceeds 5%, so it was necessary to devise such as adjusting the significance level\footnote{If you are wondering what I'm talking about, please go back to the text of basic data analysis and check. To explain it simply, the term of "5% level" was a judgment standard that let's keep the possibility of making a mistake in a single decision at 5%, but if you repeat the decision twice, it becomes a 9.75% level with $1-(0.95^2)=0.0975$. When comparing the average values of three groups, you need to compare them at three places, and repeating \keyterm{t-test} three times results in a problem of not being a level of 5%.}.


% ところが，\underline{ベイズ推定の場合はこうした問題が生じません}。なぜなら，確率を伴う判断をしていないからです。帰無仮説検定の場合は，「帰無仮説が正しいという条件のもとで」計算された検定統計量($F$値や$t$値)が，帰無仮説の条件下ではどれほど生じやすいかということを考え，その確率$p$値を計算していましたが，この$p$値はパラメータがどこにあるかといった確率ではありません。判断のための指針としての確率に過ぎないのです。一方ベイズ推定する場合，推定されたパラメータの\keytermL{事後分布}{じごぶんぷ}は，パラメータがこの辺りにあるだろう，という確率であり，生成量である差の分布もさの大きさがこれぐらいだろうという範囲，大きさに直接関わる確率です。ですから何\%区間で考えてもいいですし，その判断が間違うかどうかの確率というのは問題に含まれないのです\footnote{もっとも，すべて研究者の仮定したモデルのもとで，という話ではありますから，その仮定が間違っていたら全部意味のない数字である，ということにはなるかと思います。とはいえ，帰無仮説検定もデータが正規分布に従うという仮定で行うのですから，これについてはどっこいどっこい，というところでしょうか。}。検定の多重性の問題が生じない，というのは複雑な要因計画を考える場合には非常に助かりますね。
However, such problems do not arise in the case of \underline{Bayesian estimation}. This is because we are not making judgments with probabilities. In the case of null hypothesis testing, we calculated the test statistic ($F$ value or $t$ value), which is computed "under the condition that the null hypothesis is true", and we thought about how likely it is to occur under the condition of the null hypothesis and calculated its probability $p$ value. But this $p$ value is not a probability of where the parameter is. It is just a probability used as a guide for judgment. On the other hand, when we make a Bayesian estimate, the \keytermL{posterior distribution} of the estimated parameter is a probability that the parameter might be around here. This distribution of differences, too, is a probability directly related to the range, the size, of how much difference there might be. So you can think in terms of any percentage interval, and the probability of making a wrong judgment is not a problem\footnote{Of course, this is all under the assumption of the model proposed by the researcher, so if that assumption is wrong, it can all become meaningless numbers. However, since null hypotheses testing is also based on the assumption that the data follows a normal distribution, this might be a draw.}. The absence of problems with test multiplicity is very helpful when considering complex factorial designs.

% \section{パラメータの変形と制約}
\section{Transformation and Constraints of Parameters}
% \subsection{パラメータの変形}
\subsection{Transformation of Parameters}
% ところで，先ほどは差の分布を\verb|generated quantities|ブロックで表現しましたが，違う表現も可能です。差分も未知なるパラメータと考えて推定するのです。
By the way, we expressed the distribution of differences in the \verb|generated quantities| block earlier, but other expressions are also possible. Think of the difference as an unknown parameter and estimate it.

% 順を追って説明しましょう。まず，最初の群の平均値を$\mu_1$とします。第二の群平均$\mu_2$は，$\mu_1$と$\delta_1$だけ違う，つまり$\mu_2=\mu_1+\delta_1$と考えるのです\footnote{ここで$\delta$はデルタと読む，ギリシア文字の$d$のことです。先ほどは生成量として\verb|diff|と書いていましたが，ここでは未知のパラメータなのでギリシア文字に書き換えました。}。同様に，第三の群平均$\mu_3$は$\mu_3 = \mu_1+\delta_2$と考えます。こうすると，推定するべきパラメータは$\mu_1,\delta_1,\delta_2$であり，\verb|generated quantities|ブロックを使わなくても，直接差分パラメータを推定できるようになります。これを書いてみたのが，code:\ref{stan::21_02}になります。
Let me explain it in order. First, let's denote the mean of the first group as $\mu_1$. The mean of the second group, $\mu_2$, is different from $\mu_1$ by $\delta_1$, that is, we consider $\mu_2=\mu_1+\delta_1$\footnote{Here, $\delta$ is read as delta, which is the Greek letter for 'd'. I wrote it as \verb|diff| as a generation quantity earlier, but here it's an unknown parameter, so I rewritten it in Greek letters.}. Similarly, the mean of the third group, $\mu_3$, is considered as $\mu_3 = \mu_1+\delta_2$. In this way, the parameters to be estimated are $\mu_1,\delta_1,\delta_2$, and you can directly estimate the difference parameters without using the \verb|generated quantities| block. What I wrote is shown in code:\ref{stan::21_02}.

\begin{lstlisting}[language=C++,label={stan::21_02},caption={差分を直接推定するコード}]
...
parameters{
  real mu1;
  real delta1;
  real delta2;
  real<lower=0> sig;
}

model{
  // likelihood
  X1 ~ normal(mu1,sig);
  X2 ~ normal(mu1+delta1,sig);
  X3 ~ normal(mu1+delta2,sig);
  // prior
  mu1 ~ uniform(0,100);
  delta1 ~ normal(0,100);
  delta2 ~ normal(0,100);
  sig ~ cauchy(0,5);
}

\end{lstlisting}


% \verb|data|ブロックに違いはありませんので省略してあります。変更点として，まず\verb|parameters|ブロックで推定するべきパラメータが変わっていることを確認してください。
The \verb|data| block is omitted because there is no difference. As a change, please confirm first that the parameters to be estimated in the \verb|parameters| block have changed.
% 次に\verb|model|ブロックで，3つのデータがそれぞれ\verb|normal(mu1,sig),normal(mu1+delta1,sig),normal(mu1+delta2,sig)|から出てきているように書き変わっています。正規分布の位置パラメータをずらすことで書き換えているのです。
In the next \verb|model| block, it is rewritten as if the three pieces of data are each coming from \verb|normal(mu1,sig),normal(mu1+delta1,sig),normal(mu1+delta2,sig)|. It is being rewritten by shifting the position parameter of the normal distribution.
% 最後にその下，事前分布の設定ですが，差分$\delta_1,\delta_2$は正負どちらにも出てくる可能性があり，その大きさがわかりませんので，平均を0にした正規分布としてあります。
Finally, under that, regarding the setting of the prior distribution, the differences $\delta_1, \delta_2$ could appear on the positive or negative side, and since we don't know their size, they are recorded as normal distributions with a mean of 0.
% これで推定した結果は，先ほどの生成量を使った結果と変わりはありません。数式の展開で書き方が変わっただけで，モデルの形は同じだからです。
The estimated result with this is no different from the result using the previously mentioned production amount. The way of writing changed with the expansion of the equation, but the shape of the model is the same.

% ここで\verb|transformed parameters|ブロックの紹介をしておきましょう。code:\ref{stan::21_02}は悪くはないのですが，モデルのところ，\verb|normal(mu1+delta1,sigma)|としているのがちょっと美しくないですね。元のコードにあった，\verb|normal(mu2,sigma)|のほうがわかりやすいのに，と思った人もいると思います。まあ今の所，そこまで複雑ではないのでそれほど問題にはならないでしょうが，これからもしパラメータが複雑な構造をするようになったとき，尤度のところはもう少しスッキリ書きたいということも生じてくるでしょう。そういうときに\verb|transformed parameters|ブロックを使います。このブロックは，\verb|parameters|ブロックの次に書くもので，言葉の通りパラメータを変形(トランスフォーム)するためのものです。
Let's introduce the \verb|transformed parameters| block here. Code:\ref{stan::21_02} is not bad, but at the model part, writing it as \verb|normal(mu1+delta1,sigma)| is a bit inelegant. Some may think that it would be easier to understand if it was as \verb|normal(mu2,sigma)| in the original code. Well, it's not too complicated for now, so it's not a big issue, but if the parameters become more complex in the future, you might want to write the likelihood part a bit more neatly. That's when you use the \verb|transformed parameters| block. This block is written after the \verb|parameters| block and is, as the name suggests, used to transform parameters.

% これを使って書いたのがcode:\ref{stan::21_03}です。
The one that was written using this is code:\ref{stan::21_03}.
\begin{lstlisting}[language=C++,label={stan::21_03},caption={パラメータ変換を入れたコード}]
...
parameters{
  real mu1;
  real delta1;
  real delta2;
  real<lower=0> sig;
}

transformed parameters{
    real mu2;
    real mu3;
    mu2 = mu1 + delta1;
    mu3 = mu1 + delta2;
}

model{
  // likelihood
  X1 ~ normal(mu1,sig);
  X2 ~ normal(mu2,sig);
  X3 ~ normal(mu3,sig);
  // prior
  mu1 ~ uniform(0,100);
  delta1 ~ normal(0,100);
  delta2 ~ normal(0,100);
  sig ~ cauchy(0,5);
}

\end{lstlisting}

% ここで注目すべきは，まず\verb|transformed parameters|ブロックで\verb|mu2,mu3|という変数名が宣言され，それが式によって表現されていることです。ここでパラメータ\verb|mu1,delta1|が\verb|mu2|に変形されていることがわかります。このようにして作られたパラメータであれば，\verb|model|ブロックの中で使ってやることができます。現に，尤度のところが\verb|X2 ~ normal(mu2,sig)|のようにスッキリした形になっていますね。ただし注意すべきは，事前分布のところが変わっていない，ということです。事前分布はパラメータに対しておかれますので，ここで\verb|mu2 ~ uniform(0,100);|などと置くべきではありません。
The first thing to note here is that the variable names \verb|mu2, mu3| are declared in the \verb|transformed parameters| block, and are expressed by the formula. It can be seen here that the parameters \verb|mu1, delta1| are being transformed into \verb|mu2|. Such parameters can be used in the \verb|model| block. Indeed, the likelihood part has become neat, as in \verb|X2 ~ normal(mu2, sig)|. However, what should be noted is that the prior distribution part has not changed. The prior distribution is placed on the parameter, so it is not appropriate to put \verb|mu2 ~ uniform (0,100);| here.

% このように変換することで，コードが多少なりともスッキリしたと思います。面倒だな，と思うかもしれませんが，コードがスッキリすることは自分にとっても誰にとっても「読みやすい」ことにつながります。そして読みやすいコードは，間違い(バグ)があったときに修正しやすいのです。あまりごちゃごちゃしたコードを書くと，問題が生じたときに対応が難しくなりますから，なるべくこうした作法を活用して，コードは単純明快なものにしておくと良いでしょう。
I think the code has been tidied up to some extent by making this conversion. You might think it's a hassle, but a neat code is easier to read for you and everyone else. And easy-to-read code makes it easier to fix when there's a mistake (bug). If you write too messy code, it will be difficult to deal with when problems arise, so it's best to keep your code simple and clear by using these conventions.
% \subsection{パラメータの制約}
\subsection{Constraints of Parameters}
% さてここでこのモデルをもう一段階，発展させてみましょう。
Now, let's take this model to the next level and develop it further.
% 今のコードは第一の群，$\mu_1$を基準に，$\mu_2,\mu_3$を構成していました。別に第二，第三の群を基準にしても良かったですよね。つまり基準点は恣意的になっています。
The current code structures $\mu_2,\mu_3$ based on the first group, $\mu_1$. We could have also used the second or third group as the reference, right? In other words, the reference point has been arbitrarily set.
% それで悪いということではないのですが，$\mu_1$だけ特別扱いです。そうではない作法として，全体平均$\mu$を考えて，次のように定式化してみましょう。
It's not that it's wrong, but only $\mu_1$ is treated specially. As a different method, let's consider the overall average $\mu$ and formulate it as follows.

\begin{eqnarray}
  \mu_1 &=& \mu + \delta_1 \\
  \mu_2 &=& \mu + \delta_2 \\
  \mu_3 &=& \mu + \delta_3
\end{eqnarray}

% こうすることで，全体平均からの差分としてモデルを表現できます。
By doing this, you can express the model as a difference from the overall average.
% ここで全体平均$\mu$はどういう意味があるでしょうか。群の違いを独立変数，結果のスコアを従属変数とした線形モデルの一環として考えると，全体平均は群を通じて変化のない水平線(図\ref{fig::21_02})という関係を表していることになります。
What does the overall average $\mu$ mean here? If you think of it as part of a linear model with the differences between the groups as the independent variable and the score of the results as the dependent variable, the overall average would represent a constant line (Figure \ref{fig::21_02}) throughout the groups.

% 横一線で群ごとに違いがないのですから，いわばこれは\keytermL{帰無仮説}{きむかせつ}のモデル，すなわち\keyterm{帰無モデル}{きむもでる}{Null Model}とでもいうべきものです。
Since there is no difference between the groups along the horizontal line, this is essentially what can be called a null model, namely a null hypothesis model.

% あるいは，群ごとの平均値の違いはこの全体平均からの差分で表現されます。この群ごとの違いは，その群に入ったからこそ生じた\textbf{効果}であるとも言えます。つまり，$\delta$は効果の大きさなのです。翻ってかんがえると，全体平均は「なんの処置もしなければ理論的にこの値になるはず」という操作なしの状態，天然の状態とでもいえるでしょう。\keyterm{ランダム化比較実験}{らんだむかひかくじっけん}{Randomized Control Trial}は，被験者をランダムに割り付けることによって，平均的に見れば個別の効果が相殺し合うことを用いています。平均化することで誤差が相殺するので群ごとの比較ができますし，群の平均値が変化したということは，その群に属したことの効果が現れたと考えるわけです。帰無モデルの状態は，この個別の効果を相殺させた状態であることを表しています。
Alternatively, the difference in the mean values for each group can be expressed as a deviation from this overall average. This difference within each group can also be described as an \textbf{effect} that occurred precisely because it was included in that group. In other words, $\delta$ is the size of the effect. Conversely, the overall average can be described as the 'natural' state, or the value that should be theoretically obtained if no action is taken. A \keyterm{Randomized Control Trial} assigns participants randomly, using the principle that individual effects cancel each other out on average. By averaging, errors cancel out so comparisons can be made between groups, and a change in the average value of a group indicates the effect of belonging to that group. The state of the null model indicates a state where individual effects have been canceled out.

\begin{figure}[ht]
  \begin{center}
    \includegraphics[width=0.8\linewidth,keepaspectratio]{../images/21_modeling3/Slide21_02.png}
    \caption{全体平均は水平線}
    \label{fig::21_02}
  \end{center}
\end{figure}


% さてこのように考えて，これを数式の形，モデルの形に書き起こしていくことを考えましょう。
Now, let's consider transcribing this into the form of mathematical formulas or models, considering it in this way.
% ところがこのままだと，未知のパラメータは$\mu,\delta_1,\delta_2,\delta_3$と4つあることになります。今までは$\mu_1,\mu_2,\mu_3$，あるいは$\mu_1,\delta_1,\delta_2$の3つだったのに，です。同じ数式を書き直しただけなのに，パラメータの数が変わるのはおかしいですね。
However, as it stands, there will be four unknown parameters, $\mu,\delta_1,\delta_2,\delta_3$. Until now, it was three: $\mu_1,\mu_2,\mu_3$ or $\mu_1,\delta_1,\delta_2$. It's strange that the number of parameters changes just by rewriting the same equation.
% 実はこれには1つ，制約が欠けているのです。今回は全体平均$\mu$からの相対比較になっていますから，$\delta_1+\delta_2+\delta_3=0$という関係が隠れていることになります\footnote{もしこれが$=0$にならなければ，全体平均の計算に矛盾が生じますよね。全体平均$\mu$は，$\mu = \frac{1}{3}(\mu_1+\mu_2+\mu_3)$ですから，$\mu = \frac{1}{3}(\mu+\delta_1+\mu+\delta_2+\mu+\delta_3) =\frac{1}{3}(3\mu + \sum \delta_j) = \frac{1}{3}3\mu + \frac{1}{3}\sum \delta_j$となり，最初の項が$\frac{1}{3}3\mu = \mu$ですから，第二の項は$\frac{1}{3}\sum \delta_j = 0$でないと計算があいません！}。
In fact, there is one constraint missing from this. Since we are making a relative comparison from the overall average $\mu$, it is implicit that there is a relationship $\delta_1+\delta_2+\delta_3=0$\footnote{Otherwise, there would be a contradiction in the calculation of the overall average, right? The overall average $\mu$ is $\mu = \frac{1}{3}(\mu_1+\mu_2+\mu_3)$, so $\mu = \frac{1}{3}(\mu+\delta_1+\mu+\delta_2+\mu+\delta_3) =\frac{1}{3}(3\mu + \sum \delta_j) = \frac{1}{3}3\mu + \frac{1}{3}\sum \delta_j$. Here, the first term is $\frac{1}{3}3\mu = \mu$, so the second term has to be $\frac{1}{3}\sum \delta_j = 0$ for the calculation to match!}.
% そこで，この制約をモデルの中に書き込まなければなりません。$\delta_1+\delta_2+\delta_3=0$という式から，変数を移項することでこれを表現します。すなわち$\delta_3 = 0 - (\delta_1+\delta_2)$のようにします。
Therefore, you must write this constraint into the model. This is expressed by rearranging the variables from the equation $\delta_1+\delta_2+\delta_3=0$. That is, it's like $\delta_3 = 0 - (\delta_1+\delta_2)$.

% これを使って書いたのがcode:\ref{stan::21_04}です。
I used this to write it: code:\ref{stan::21_04}.
\begin{lstlisting}[language=C++,label={stan::21_04},caption={制約を入れたコード}]
...
parameters{
  real gm;
  real delta1;
  real delta2;
  real<lower=0> sig;
}

transformed parameters{
    real delta3;
    real mu1;
    real mu2;
    real mu3;
    delta3 = 0 - (delta1 + delta2);
    mu1 = gm + delta1;
    mu2 = gm + delta2;
    mu3 = gm + delta3;
}
...
\end{lstlisting}


% このようにすることで，先ほどと等価なモデルを別の形で書き表すことができるようになりました。
By doing so, we have been able to express the same model as before in a different form.

% ここまでをまとめておきます。三群の平均値差を求めるモデルを，3つの方法で書いてみました。
Here's a summary so far. I tried to write a model that calculates the difference in averages of three groups in three different ways.
\begin{enumerate}
  \item 3つの群平均$\mu_1,\mu_2,\mu_3$を個別に推定するモデル
  \item ある群とそこからの差分，$\mu_1,\delta_1,\delta_2$を推定するモデル
  \item 全体平均とそこからの差分，$\mu,\delta_1,\delta_2,\delta_3$を推定するモデル。ただし$\sum \delta = 0$という制約をつける
\end{enumerate}

% この方法はいずれも同じ推定結果を返します。第一のモデルの生成量から差分を計算すれば，$\delta_1,\delta_2$を計算したことと同じです。
Both methods return the same estimated results. Computing the difference from the generation amount of the first model is equivalent to computing $\delta_1$ and $\delta_2$.

% 同じことを異なる方法で実装するのは，生産的でないように思うかもしれません。しかし第3の形式にしておくと，モデルの一般化が容易になります。
Implementing the same thing in different ways may seem unproductive. However, keeping it in a third form makes it easier to generalize the model.
% すなわち，多群の平均値差の比較をするときに，群の数が変わるたびにコードを書き換える手間がなくなるのです！
That is, when comparing the differences in averages among many groups, you will no longer have to rewrite the code every time the number of groups changes!
% \section{モデルの洗練}
\section{Refinement of the Model}
% ここまで，二群，三群の平均値差を求めるモデルを書いてきました。書き方はわかっているので，比較する群が四群，五群と増えても書き足していくだけでなんとかなりそうです。
Up to now, we have been writing a model to find the difference in mean values between two and three groups. Since we know how to write it, even if the number of groups we compare increases to four or five, it seems manageable just by adding to it.

% しかし，たとえば四群に増えたときに$\delta_1,\delta_2,\delta_3$,五群に増えたときに$...\delta_3,\delta_4$と手作業で増やしていると，N群までふえたときは$...\delta_{N-2},\delta_{N_1}$と何行も書き足していかなければなりません。しかもその度にコンパイルしているようでは，とても面倒なことはすぐに想像がつきます。
However, for example, when it increases to four groups, $\delta_1,\delta_2,\delta_3$, and when it increases to five groups, $...\delta_3,\delta_4$, if you are manually increasing it, when it increases to N groups, you have to write many lines like $...\delta_{N-2},\delta_{N-1}$. Moreover, if you are compiling each time, it is easy to imagine that it will be very troublesome.

% そこで，群の数が変わっても同じコードで対応できるように，一般化を目指してコードを改良しましょう。
Let's improve the code with the aim of generalizing it so that it can handle different numbers of groups with the same code.
% そのためには，何群の比較なのかを示す群の数を，\verb|data|ブロックで外部から取り込むようにすれば良いでしょう。
To do that, you should import the number of groups indicating what group comparison it is from outside in the \verb|data| block.
% その変数を使って差分の数，効果の数を計算してやれば良いのです。たとえば比較する群の数を\verb|G|とすると，全体平均$\mu$と，$G-1$個の差分をパラメータとして推定すれば良いことになります。
You should calculate the number of differences and the number of effects using that variable. For example, if the number of groups to be compared is \verb|G|, it means that you should estimate the overall average $\mu$ and $G-1$ differences as parameters.

% その考え方を使って書いたコードは，たとえばcode:\ref{stan::21_05}のようになります。
The code written using that way of thinking would be, for example, like code:\ref{stan::21_05}.

\begin{lstlisting}[language=C++,label={stan::21_05},caption={群の数を一般化したコード}]
data{
  int<lower=0> Lv;              // 水準数
  int<lower=0> N;               // 各群のサンプルサイズ       
  array[Lv,N] real X;           // 変数の値
}

parameters{
  real gm;                      // 全体平均
  array[Lv-1] real raw_delta;   // 全体からの差。水準数マイナス1個
  real<lower=0> sig;            // 誤差の分散
}

transformed parameters{
  array[Lv] real delta;           // 差の大きさを作り直す
  array[Lv] real mu;              // 再構成される群ごとの平均
  
  for(i in 1:(Lv-1)){
    delta[i] = raw_delta[i];    // ほとんどコピー
  }
  
  delta[Lv] = 0 - sum(raw_delta); // 総和が0になるように最後だけ書き換える

  for(i in 1:Lv){
    mu[i] = gm + delta[i];        // 群ごとに再構成
  }
  
}

model{
  // Likelihood
  for(l in 1:Lv){
      for(i in 1:N){
          X[l,i] ~ normal(mu[l],sig);
          
      }
  }
    
  // Prior
  gm ~ uniform(0,100);
  raw_delta ~ normal(0,100);
  sig ~ cauchy(0,5);
}

\end{lstlisting}


% \paragraph{コード解説}
\paragraph{Code Explanation}
\begin{description}
	\item[dataブロック] 群の数\verb|Lv|，各群のサンプルサイズ\verb|N|,各データ点\verb|X|を入力します。ここで\verb|X|が\textbf{二元配列}になっていることに注目してください。\verb|X[1,1]|は第一群の第1番目のデータ，\verb|X[1,2]|は第一群の第2番目のデータ，\verb|X[2,1]|は第二群の第1番目のデータ，同様に\verb|X[i,j]|は第i群の第j番目のデータを表すことになります。この形式は二元配列(2次元の変数セット)になっていると言います\footnote{行列でもいいじゃないか，と思われるかもしれません。それでも結構です。Stanは行列の時，\verb|real|ではなく\verb|matrix|で宣言することができます。しかしこのように\verb|real|型にしておいて，後ろのカッコで配列の次元を表現すると，何次元でも拡張することができるので今回はそのようにしました。}。宣言の時に，上で取り込んだ変数\verb|Lv|や\verb|N|を使うことができます。最大\verb|Lv|個の群，各群\verb|N|人のデータがあるという前提です。
	\item[parametersブロック] 全体平均\verb|gm|と，差分，誤差のSDをパラメータとしています。差分はここでは\verb|raw_delta|とし，上で宣言した\verb|Lv|をつかって\verb|Lv-1|個の要素を持つ配列にしています。\verb|raw_delta|って変な名前，と思うかもしれませんが，この後これを変換して$\delta$を作りますので，作る前の生(raw)の$\delta$という名前にしてみました\footnote{変数名は任意ですので，著者の命名法がわかりにくいなあと思ったら，自分で好きな変数名にしてもらって構いませんよ。その場合は，以後すべての変数を読み替えてくださいね。}。
	\item[transformed parametersブロック] 以後の話をわかりやすくするために，各群の平均\verb|mu|を構成することにします。\verb|mu|の数は水準数と同じ\verb|Lv|です。そして第i群の平均\verb|mu[i]|を作るために，\verb|mu[i]=gm+delta[i]|という書き方をしたいので，水準数と同じ数だけ\verb|delta[i]|を作りたいと思います。もとの\verb|raw_delta|が1,2,...,Lv-1個ありますので，\verb|delta|もLv-1番目まではそれと同じものをコピーします(\verb|for|文で繰り返し代入していってます。)。そして最後にLv番目の要素を，0からこれまでの要素をすべて足した(\verb|sum(raw_delta)|)ものを引くことで作っています。ちなみに\verb|sum|はstanの持っている総和の関数で，配列要素のすべてを足し合わせるというものです。こうして\verb|delta|がLv個できあがりましたので，各水準の平均値を\verb|mu[i]=gm+delta[i]|という数式で一般的に書くことができるようになりました。
	\item[modelブロック] 尤度のところに，群それぞれについて巡回する\verb|for|文，その中で1,2,3...N人まで巡回する\verb|for|文を組み込んで二重にぐるぐると回しています。事前分布は設定の通りです。
\end{description}


% このコードを使って，実際に推定してみましょう。データは表\ref{tbl::21_01}のものを使います。
Let's try to estimate using this code. We will use the data from Table \ref{tbl::21_01}.
\begin{table}[ht]
  \begin{center}
    \caption{独立した三群から得られたスコア}
    \begin{tabular}{l|llll}\hline
      groupA & 6 & 6 & 5 & 5 \\
      groupB & 7 & 4 & 7 & 6 \\
      groupC & 4 & 3 & 4 & 6 \\\hline
    \end{tabular}
    \label{tbl::21_01}
  \end{center}
\end{table}


% 推定してみましょう(Rコードは\ref{code:21_01},結果は\ref{MCMCout:out21_01})。
Let's try to estimate (R code is \ref{code:21_01}, the result is \ref{MCMCout:out21_01}).
% 各群の平均値，全体平均からの偏差などが，事後分布として出力されているのがわかると思います。これをみて，たとえば差の事後分布の50\%確信区間でみると$\delta_2,\delta_3$はその区間に0を挟んでないので，一定の効果はあるだろうなどと判断できるわけです。もちろん生成量をつかって，ある程度の差がある確率がどのくらいとか，優越率などデータレベルの仮説を立てることも可能です。
I think you can see that things like the mean values of each group and deviations from the overall average are output as a posterior distribution. Looking at this, for example, if you look at the 50% credible interval of the posterior distribution of the difference, $\delta_2, \delta_3$ does not include 0 in that interval, so you can judge that there is a certain effect. Of course, using the amount of production, you can also hypothesize the data level such as how much probability there is a certain difference or dominance ratio.

\begin{lstlisting}[language=c++,label={code:21_01},caption={三群の平均値の比較}]
Example <- matrix(c(6, 6, 5, 5, 7, 4, 7, 6, 4, 3, 4, 6), 
                  ncol = 4, byrow = T)
dataSet <- list(Lv = 3, N = 4, X = Example)
\end{lstlisting}


\begin{MCMCout}{}{out21_01}
\begin{verbatim}
# A tibble: 10 × 7
   name               EAP       MED       MAP        SD       L95       U95
   <chr>        <num:.3!> <num:.3!> <num:.3!> <num:.3!> <num:.3!> <num:.3!>
 1 delta[1]         0.257     0.256     0.227     0.569    -0.880     1.394
 2 delta[2]         0.749     0.754     0.804     0.566    -0.387     1.878
 3 delta[3]        -1.005    -1.009    -1.114     0.576    -2.158     0.162
 4 gm               5.251     5.251     5.242     0.400     4.452     6.039
 5 mu[1]            5.507     5.507     5.562     0.699     4.114     6.902
 6 mu[2]            5.999     6.001     5.986     0.697     4.603     7.386
 7 mu[3]            4.246     4.238     4.182     0.695     2.860     5.647
 8 raw_delta[1]     0.257     0.256     0.227     0.569    -0.880     1.394
 9 raw_delta[2]     0.749     0.754     0.804     0.566    -0.387     1.878
10 sig              1.339     1.261     1.134     0.391     0.823     2.301
\end{verbatim}
\end{MCMCout}


% \subsection{データサイズの一般化}
\subsection{Generalization of Data Size}
\label{tidyDataForm}
% さてこのモデルは，比較する群の数が増えても外部からデータとして群の数を与えることができますので，毎回コンパイルする必要がありません。やったね！
Now, with this model, there's no need to compile every time because it allows you to provide the number of groups as data from outside, even if the number of groups to compare increases. Well done!
% これで完璧，と言いたいところですが，1つ気になるのはサンプルサイズです。各群のサンプルサイズ$N$をデータとして与えるようになっていますが，たとえばA群は10人，B群は12人，C群は14人と言ったように，群ごとにサンプルサイズが異なるとこのコードでは対応できません。実際に実験をやる場合，各群のサイズを整えて人を集めたいところですが，調査研究などでは群ごとに同じ人数を与えるような調整ができないこともあり，そういう意味ではこのコードではうまく対応できそうにありません。
I would like to say it's perfect but there is one concern about the sample size. Each group's sample size $N$ is given as data. However, this code cannot handle cases where the sample size is different for each group, for example, 10 for group A, 12 for group B, and 14 for group C. In actual experiments, you would want to match the size of each group and gather people, but in surveys and studies, it may not be possible to adjust to give the same number of people in each group. In this sense, this code may not be able to handle it properly.

% そこで，群の人数が異なる場合でも対応できるように，さらにこのコードを拡張していきたいと思います。そのためにまず，データを\keyterm{整然データ}{せいぜんでーた}{tidy data}の形に整形しましょう。
So, in order to handle cases where the number of people in the group differs, I would like to further extend this code. To do that, let's first shape the data into a form of \keyterm{tidy data}.

% 先ほどのサンプルデータ(表\ref{tbl::21_01})は，$3\times4$の長方形をしていました。これはデータのサイズが整っているからできることで，サイズが変わってしまうと表の中に欠損ができてしまうことになります。
The sample data (Table\ref{tbl::21_01}) that we had earlier was in the shape of a $3\times4$ rectangle. This is possible because the size of the data is in order; if the size changes, missing parts would appear in the table.
% また，たとえばB群の3番目の人の値を見る時，われわれは2行3列目のデータにさっと目が行きますが，機械的には群の情報が行名に，人の群内整理番号が列名にあるので，1つのデータを見る時に行・列それぞれを参照することになります(図\ref{fig::21_03}左)。ここで「どの群か」「群の中の何番目の人か」という情報はいずれも変数で，人によって変わるものですから，データの持つ情報の1つなのです。データの持つ情報，変数なのに参照先が変わっていることになります。
For example, when we look at the value of the third person in group B, our eyes quickly go to the data in row 2, column 3, but mechanically, the group information is in the row name and the person's internal group number is in the column name. So when we look at a piece of data, we have to refer to both the row and the column (Figure\ref{fig::21_03}left). Here, both the information of "which group" and "which person in the group" are variables, and they vary from person to person, so they are one of the pieces of information that the data holds. Even though they are pieces of information held by the data, or variables, their references have changed.

% このデータの情報をまったく損なうことなく，縦長に並べ直したのが図\ref{fig::21_03}の右側です。これは1つの行がそれぞれの観測に対応しています。このような形式にすると，何行目かということを指定するだけで，どの群の何番目の人なのかという情報が手に入ります。またデータの中に欠損があっても，整然データの形にする時にこれは観測に含まれないので，データの一部が欠けることがないという利点があります。
The right side of Figure \ref{fig::21_03} in the data re-arranges it lengthwise without losing any information. This corresponds to the observation for each row. In such a format, we can find the information of which group belongs to whom by just specifying the row number. Furthermore, even if there are missing values in the data, this is not included in the observation when tidying up the data, so there is an advantage that no part of the data is missing.
\begin{figure}[ht]
  \begin{center}
    \includegraphics[width=0.8\linewidth,keepaspectratio]{../images/21_modeling3/Slide21_03.png}
    \caption{一般的なデータと整然データ}
    \label{fig::21_03}
  \end{center}
\end{figure}


% このようなデータ形式にして，これを与える形でデータ分析をするとデータサイズに依存しないコードが書けます。具体的にはcode:\ref{stan::21_06}のようにします。
Use such a data format, and if you perform data analysis in this way, you can write a code that does not depend on the data size. Specifically, do it as in code:\ref{stan::21_06}.
\begin{lstlisting}[language=C++,label={stan::21_06},caption={整然データに対応したコード}]
data{
  int<lower=0> Lv;                             // 水準数
  int<lower=0> L;                              // データ数
  array[L] int<lower=0,upper=Lv> idx;          // データのID
  array[L] real X;                             // 変数の値
}
...(中略)...
model{
  // Likelihood
  for(l in 1:L){
    X[l] ~ normal(mu[idx[l]],sig);
  }
    
  // Prior
  gm ~ uniform(0,100);
  raw_effect ~ uniform(-100,100);
  sig ~ uniform(0,100000);
}
\end{lstlisting}


% \verb|parameters|ブロックと\verb|transformed parameters|ブロックはcode:\ref{stan::21_05}と同じなので省略しました。
The \verb|parameters| block and the \verb|transformed parameters| block have been omitted because they are the same as in code:\ref{stan::21_05}.
% まず\verb|data|ブロックをみてください。群の数を\verb|Lv|で入力するのはそのままですが，次にデータの総数\verb|L|を撮るようにしています。図\ref{fig::21_03}の例で言えば，C群3列目のデータがありませんから，データ長は全部で$L=11$になります。
First, please look at the \verb|data| block. Entering the number of groups in \verb|Lv| is as is, but next, we are taking the total number of data \verb|L|. For instance, in the example of figure \ref{fig::21_03}, there is no data in the third column of group C, so the total data length is $L=11$.

% 次にL行のデータがどの群に属するのかを指示するインデックス変数\verb|idx|をL個用意しています。L行目のデータが第何群に属しているのかの数字が入ります。たとえば1行目はA群なので1，5行目(L=5)はB群なので2，と言った数字が入るようになっています。
Next, we have prepared the index variable \verb|idx|, which indicates to which group the L rows of data belong. It will contain the number of which group the data of the Lth row belongs to. For example, since the first line is group A, it will be 1, and since the fifth line (L=5) is group B, it will be 2, and so forth.
% さいごにデータの値そのものである\verb|X|も，データ長Lと同じだけの1次元配列を用意してやります。
Finally, prepare a one-dimensional array of the same length L as the data for the value of the data itself, \verb|X|.

% 技術的に注目すべきは\verb|model|ブロックの尤度のところで，\verb|for|文がデータ長Lの反復をしているだけです。つまり1,2,3...,L行目のデータを順に参照していくのです。
Technically notable is the likelihood part of the \verb|model| block, where the \verb|for| statement is just iterating over the data length L. In other words, it refers to the data lines of 1, 2, 3..., L in sequence.
% そして$l$行目のデータがどの群に属するかは，変数\verb|idx[l]|が指し示してくれますから，\verb|mu[idx[l]]|としてやることで指定できていることになります。\verb|idx[l]|は，$l$行目のデータが属している群の数字が代入されていますから，たとえば1行目(l=1)の場合は\verb|mu[idx[1]] = mu[1]|としていることと同じ，5行目の場合は\verb|mu[idx[5]]=mu[2]|としていることと同じ，ということになります。ここでは\underline{入れ子になった参照}が行われています。ちょっとテクニカルですが，この技術を使うと表現力も広がりますので，仕組みをしっかり理解しておいてください。
The variable \verb|idx[l]| indicates which group the data in the $l$-th row belongs to, so you can specify it by doing \verb|mu[idx[l]]|. Since \verb|idx[l]| is assigned with the number of the group that the data in the $l$-th row belongs to, for example, if it is the first row (l=1), it is equivalent to doing \verb|mu[idx[1]]=mu[1]|, and if it is the fifth row, it is equivalent to doing \verb|mu[idx[5]]=mu[2]|. Here, a \underline {nested reference} is made. It's a bit technical, but using this technique expands your expressive power, so please make sure you understand the mechanism.

% \section{パラメータリカバリ}
\section{Parameter Recovery}
% さあ，最後にパラメータリカバリをして，このコードが正しく推定できるか，あるいはどの程度のサンプルサイズがあればどの程度の精度で推定できるのかを検証しておきましょう。
Now, let's finally do parameter recovery, and verify whether this code can estimate correctly, or to what extent of accuracy can it estimate with what size of sample.

% 私たちはデータがどのように造られるか，というそのメカニズムの方からアプローチしているわけです。これはリバースエンジニアリングと呼ばれる考え方ですね。そして今から仮想データを作ろうという時も，そのデータ生成モデルをそのまま利用すればいいのです。仮想データの生成は，データ生成メカニズムをリバースエンジニアリングで明らかにしていくアプローチのちょうど正反対，リバース・リバース・エンジニアリングです。
We approach from the perspective of how data is created. This concept is called reverse engineering. And when we are about to create virtual data, we can simply use the data generation model. The generation of virtual data is the opposite of the reverse engineering approach to elucidating the data generation mechanism, which is reverse-reverse-engineering.

% 具体的には，たとえば次のようなコード(code:\ref{code:21_02})でデータを作ることができるでしょう。
Specifically, for example, you can create data with the following code (code:\ref{code:21_02}).
\begin{lstlisting}[language=c++,label={code:21_02},caption={リバースエンジニアリング}]
N <- 100
Lv <- 5
gm <- 50
sig <- 3

raw_effect <- runif(Lv - 1, -10, 10)
effect <- c(raw_effect, 0 - sum(raw_effect))
mu <- gm + effect

X <- rnorm(N * Lv, mu, sig)
dat <- data.frame(
  Idx = rep(1:Lv, N * Lv),
  value = X
)

modelC <- cmdstanr::cmdstan_model("BetweenAnova2.stan")
dataSet <- list(Lv = Lv, L = NROW(dat), idx = dat$Idx, X = dat$value)
\end{lstlisting}

% \paragraph{コード解説}
\paragraph{Code Explanation}
\begin{description}
	\item[1行目] 各群のサンプルサイズNを設定します。各群共通のサイズになります。
	\item[2行目] 水準数です。ここでは5群の平均値の比較をすることにしました。
	\item[3行目] 全体平均の設定です。事前分布に[0,100]の一様分布を置いてますから，真ん中ぐらいにしてあげました。
	\item[4行目] 誤差の散らばりを設定します。
	\item[6行目] 水準数$-1$の効果を設定します。手入力でもいいのですが面倒なので，-10から10までの範囲で一様乱数によって生成しました\footnote{一様乱数の関数は，Rでは\texttt{runif}です。}。
	\item[7行目] 水準数Lvの効果にするため，先ほど作った水準数$-1$の\verb|raw_effect|に\verb|0-sum(raw_effect)|の計算結果をつけくわえています。関数\texttt{c}は結合させるcombineという意味です。
	\item[8行目] 各群の平均値です。全体平均に効果を足しています。全体平均は1つの数字，効果は水準数の要素を持つベクトルですから，計算ができないように思えますが，このような場合Rは自動的にサイズ合わせのため値の再利用を行います。つまり，$50+(\delta_1,\delta_2,\delta_3,\delta_4,\delta_5)$を$(50,50,50,50,50)+(\delta_1,\delta_2,\delta_3,\delta_4,\delta_5)$と解釈してくれるのです。
	\item[10行目] 水準数$\times$各群のサイズの乱数を発生させています。乱数は正規分布によるもので，平均やSDは上で指定した通りです。ここでも値の再利用が行われていて，ベクトル\verb|mu|は$\mu_1,\mu_2,\mu_3,\mu_4,\mu_5,\mu_1,\mu_2,....$とくり返して当てはめられていきます。
	\item[11-14行目] 数値を\verb|data.frame|型にくみ上げていってます。\verb|rep|関数は，指定した回数だけ繰り返すもので，ここでは\verb|1:Lv|すなわち1,2,3,4,5という数字をデータサイズ分繰り返していることになります。
	\item[16行目] ここでは\verb|cmdstanr|パッケージを使ってコンパイルしています。
	\item[17行目] 推定に使うデータセットです。データ長は\verb|data.frame|の長さを返す関数\verb|NROW|を使って与えています。インデックス変数は\verb|dat|オブジェクトの\verb|Idx|変数，データは同じオブジェクトの\verb|value|変数を指定しています。
\end{description}


% いかがでしょうか。作ったデータ，オブジェクト，ベクトルなどの意味がわからない場合は，その都度Rでオブジェクト名を入力し，何が格納されているかを確認しながら進めると良いでしょう。
How about that? If you don't understand the meaning of the data, objects, vectors, etc., that you've created, it's a good idea to type in the object name in R each time and proceed while checking what is stored.
% ここで作った仮想データを，先ほどのStanオブジェクトに与えて乱数を生成し，パラメータリカバリがどの程度の精度でできるかを考えてみてください。サンプルサイズや誤差の大きさなどを変えながら確認してみると良いでしょう。
Please consider creating random numbers by giving the virtual data created here to the Stan object from earlier, and see how accurately parameter recovery can be achieved. It would be good to check this by changing things such as the sample size and the magnitude of the error.
% \section{課題}
\section{Assignment}
% 今回は，基本課題と発展課題の2つを用意します。基本課題は必須ですが，応用課題は提出者に加点されるだけで必須ではありません(提出しなかったからと言って減点されることもありません)。
This time, we are preparing two tasks: a basic task and an advanced task. The basic task is mandatory, but the advanced task is not mandatory, but extra points will be awarded to those who submit it (you will not be penalized for not submitting it).


% \paragraph{基本；整然データに対応した多群比較のコードを書く}最後に紹介した，整然データに対応した多群比較のコードを書き，パラメータリカバリのコードとデータを使って正しく動くかどうか確認してください。作ったStanファイルやRコードを提出してください。
\paragraph{Basic: Write a code for multi-group comparisons corresponding to neat data} Finally, please write the code for multi-group comparisons corresponding to the neat data we introduced last, and verify whether it works correctly using the parameter recovery code and data. Please submit the Stan files and R codes you created.
% \paragraph{発展；自分のデータを使って分析する}基礎実験1や基礎実験2など，他の授業でとったいくつかの群の平均値を比較するようなデータを用意し，今回のモデルを適用して平均値の比較を行なってください。とくにそのようなデータセットを持っていない場合は，こちらから提供するサンプルデータを使ってください。
\paragraph{Development; Analyzing Using Your Own Data} Please prepare data that compares the average values of several groups obtained in other classes, such as Basic Experiment 1 and Basic Experiment 2, and apply this model to compare average values. Especially if you do not have such a data set, please use the sample data provided by us.

\end{document}
